% © 2026 Ing. David Jaroš — CC BY-NC-ND 4.0
%
% This work is licensed under a Creative Commons Attribution-NonCommercial-NoDerivatives
% 4.0 International License (CC BY-NC-ND 4.0).
%
% See LICENSE.md for full license text.
%
% File: canonical/alpha/alpha_equation_matrix.tex
% Purpose: Complete equation chain for all three active alpha routes.
%   Route A (modular_hecke_route): derive B from Hecke/modular structure
%   Route B (theta_spectral_v_eff_route): V_eff prime attractor (PRIMARY)
%   Route C (electroweak_or_gut_bridge): GUT bridge sin^2theta_W(GUT)=3/8
%
% Status: 2026-04-29

\ifdefined\INCLUDEMODE
\else
  \documentclass[12pt,a4paper]{article}
  \usepackage[a4paper, margin=2.5cm]{geometry}
  \usepackage{amsmath, amssymb, amsthm}
  \usepackage{mathtools}
  \usepackage{hyperref}
  \usepackage{xcolor}
  \usepackage{booktabs}
  \usepackage{tcolorbox}
  \tcbuselibrary{skins,breakable}
  \usepackage{mdframed}

  \newtheorem{theorem}{Theorem}[section]
  \newtheorem{lemma}[theorem]{Lemma}
  \newtheorem{proposition}[theorem]{Proposition}
  \newtheorem{corollary}[theorem]{Corollary}
  \theoremstyle{definition}
  \newtheorem{definition}[theorem]{Definition}
  \theoremstyle{remark}
  \newtheorem{remark}[theorem]{Remark}

  \newcommand{\BB}{\mathbb{B}}
  \newcommand{\CC}{\mathbb{C}}
  \newcommand{\RR}{\mathbb{R}}
  \newcommand{\HH}{\mathbb{H}}
  \newcommand{\ZZ}{\mathbb{Z}}
  \newcommand{\Tr}{\mathrm{Tr}}
  \newcommand{\re}{\mathrm{Re}}

  \newcommand{\PROVED}{\textbf{[Proved]}}
  \newcommand{\OPEN}{\textbf{[Open]}}
  \newcommand{\BLOCKED}{\textbf{[Blocked]}}
  \newcommand{\Lone}{\textbf{[L1]}}
  \newcommand{\Lzero}{\textbf{[L0]}}
  \newcommand{\Ltwo}{\textbf{[L2]}}

  \title{\textbf{Alpha Derivation: Equation Matrix for Three Active Routes}}
  \author{Ing.\ David Jaro\v{s}}
  \date{April 2026}

  \begin{document}
  \maketitle
\fi

%──────────────────────────────────────────────────────────────────────────────
\section{Purpose and Active Routes}
\label{sec:alpha_eq_purpose}
%──────────────────────────────────────────────────────────────────────────────

This document presents the complete equation chain for the three
currently active alpha derivation routes, with explicit identification of
proved steps, open gaps, and hard blockers.
All three routes share the same proved foundation and diverge at the
point of determining the effective potential coefficient~$B$.

\begin{tcolorbox}[colback=blue!3!white,colframe=blue!50!black,title=Active Routes]
\begin{tabular}{lll}
\toprule
\textbf{Route} & \textbf{Name} & \textbf{Target} \\
\midrule
A & Modular-Hecke & Derive $B$ from $\mu(\Gamma_0(n^*))/3$ via $S[\Theta]$ \\
B & V$_{\mathrm{eff}}$ spectral (PRIMARY) & $n^*(B_{\mathrm{phenom}}) = 137$; attack Gap G137-B \\
C & EW/GUT bridge & $\sin^2\theta_W(\mathrm{GUT}) = 3/8$ via SU(5) embedding \\
\bottomrule
\end{tabular}
\end{tcolorbox}

\textbf{Kill condition} (applies to all routes): a route is killed if it requires
a free parameter fitted to reproduce $\alpha$ or 137.

%──────────────────────────────────────────────────────────────────────────────
\section{Shared Foundation (All Routes)}
\label{sec:alpha_foundation}
%──────────────────────────────────────────────────────────────────────────────

All three routes build on the following proved results.
These are not inputs to be questioned in this document --- they are
the starting platform.

\subsection{N$_{\mathrm{eff}}$ = 12 from Algebra}

\begin{theorem}[Mode count, \PROVED\ \Lzero]
\label{thm:neff}
The biquaternion algebra $\BB = \CC\otimes_\RR\HH$ admits exactly
$N_{\mathrm{eff}} = 12$ independent real bosonic modes in its kinetic term:
\[
  N_{\mathrm{eff}}
  \;=\;
  \underbrace{8}_{\dim_\RR\BB}
  \;+\;
  \underbrace{3}_{\dim_\RR(\mathrm{Im}\,\HH)}
  \;+\;
  \underbrace{1}_{\mathrm{phase}}
  \;=\; 12.
\]
This count is reproduced by five independent derivation routes;
see \texttt{canonical/n\_eff/}.
\end{theorem}

\subsection{V$_{\mathrm{eff}}$ Structure}

\begin{theorem}[Effective potential, \PROVED\ \Lone]
\label{thm:veff}
The one-loop effective potential for winding modes on $S^1_\psi$ takes the form
\[
  V_{\mathrm{eff}}(n)
  \;=\;
  n^2 - B \cdot n \cdot \ln n,
  \qquad n \in \ZZ_{>0},
\]
where $B$ is a positive real coefficient determined by the one-loop integral
of $S[\Theta]$ around the $\psi$-circle.
Source: \texttt{canonical/alpha/alpha\_best\_route.tex §2}.
\end{theorem}

\subsection{V$_{\mathrm{eff}}$ Minimum}

\begin{theorem}[Prime attractor, \PROVED\ \Lone\ (given $B$)]
\label{thm:prime_attractor}
The continuous minimum of $V_{\mathrm{eff}}(n) = n^2 - B\,n\ln n$ satisfies the
\textbf{transcendental equation}
\begin{equation}
  2n^* = B\,(\ln n^* + 1).
  \label{eq:nstar_transcendental}
\end{equation}
This has no elementary closed form; it is solved numerically.
Selected values: $n^*(8\pi)\approx 65$; $n^*(41.57)\approx 120$; $n^*(46.28)\approx 137$.
For the phenomenological value $B = B_{\mathrm{phenom}} \approx 46.298$,
numerical solution gives $n^*_{\mathrm{continuous}} \approx 137.05$, and the
prime minimiser is
\[
  n^*_{\mathrm{prime}} = 137.
\]
This is proved conditional on $B = B_{\mathrm{phenom}}$.
Source: \texttt{canonical/alpha/alpha\_best\_route.tex~\S7};
see also \texttt{canonical/alpha/veff\_corrected\_statement.tex}.

\textbf{Note}: The formula $n^*(B) = e^{(2-B)/B}\cdot e$ previously stated here
was incorrect (it equals $e^{2/B}\approx 1.04$ for $B=46.3$) and has been
removed. See \texttt{reports/alpha\_veff\_sanity\_audit.md~\S6}.
\end{theorem}

\subsection{One-Loop Baseline}

\begin{theorem}[Baseline coefficient, \PROVED\ \Lone]
\label{thm:b0}
The one-loop contribution from the kinetic sector $S_{\mathrm{kin}}[\Theta]$
on $S^1_\psi$ gives the baseline coefficient:
\[
  B_0 \;=\; 8\pi \;\approx\; 25.13.
\]
With $B = B_0$, the prime attractor gives $n^*(B_0) \approx 65$, not 137.
The gap is a factor $B_{\mathrm{phenom}}/B_0 \approx 1.84$.
Source: \texttt{canonical/t\_munu/}.
\end{theorem}

\medskip
\textbf{Gap G137-B} (shared by routes A and B):
Derive $B = B_{\mathrm{phenom}} \approx 46.298$ from $S[\Theta]$
without using $\alpha$ as input.
The missing factor is approximately $1.84 = B_{\mathrm{phenom}} / B_0$.

%──────────────────────────────────────────────────────────────────────────────
\section{Route A: Modular-Hecke Route}
\label{sec:route_a}
%──────────────────────────────────────────────────────────────────────────────

\subsection{Central Claim}

\begin{tcolorbox}[colback=green!4!white,colframe=green!60!black,
  title=Route A Target]
Derive the B-coefficient from the modular structure of $S[\Theta]$:
\[
  B \;=\; \frac{\mu(\Gamma_0(137))}{3}
        \;\approx\; \frac{138.00}{3}
        \;=\; 46.00
        \;\approx\; B_{\mathrm{phenom}} \quad (0.64\%\text{ discrepancy}).
\]
Here $\mu(\Gamma_0(n))$ is the index of $\Gamma_0(n)$ in $\mathrm{SL}_2(\ZZ)$,
which equals $n\prod_{p|n}(1+p^{-1})$.
\end{tcolorbox}

\subsection{Equation Chain}

\paragraph{Step A1: Identify $S[\Theta]$ as a modular object.}
The action $S[\Theta]$ on the torus $T^2 = S^1_t \times S^1_\psi$,
evaluated for winding modes $\Theta_n(t,\psi) = e^{2\pi i n\psi/R_\psi} \otimes \theta_n(t)$,
takes the form:
\[
  S_n[\Theta]
  \;=\;
  \int_0^{2\pi R_t}\!\mathrm{d}t
  \int_0^{2\pi R_\psi}\!\mathrm{d}\psi\;
  \re\bigl[\Tr\bigl(\nabla^\dagger\nabla\Theta_n\bigr)\bigr]
  \;=\;
  \frac{4\pi^2 n^2}{R_\psi^2}\cdot\mathcal{F}(R_t,R_\psi).
\]

\paragraph{Step A2: Evaluate the one-loop determinant.}
The one-loop path integral over fluctuations $\delta\Theta$ around $\Theta_n$
gives a functional determinant involving the heat kernel on $S^1_\psi$.
The heat kernel trace is:
\[
  K(t_E;R_\psi)
  \;=\;
  \sum_{n=-\infty}^{+\infty} e^{-n^2 t_E / R_\psi^2}
  \;=\;
  \vartheta_3\!\left(0\,\big|\,\frac{it_E}{\pi R_\psi^2}\right).
\]

\paragraph{Step A3: Extract B from the Kac-Moody level.}
The WZW boundary term in $S[\Theta]$ on $S^1_\psi$ introduces a Kac-Moody level
$k$ modifying the heat kernel coefficient.
\begin{mdframed}[backgroundcolor=yellow!15]
\textbf{Gap A-1} \OPEN: Compute the Kac-Moody level $k$ from the boundary
structure of $S[\Theta]$ on $\partial(S^1_\psi \times M^4)$.\\
\textit{Expected}: $k = 1$ from the linear $\psi$-coupling in $\nabla^\dagger$.\\
\textit{Impact}: If $k = 1$, the one-loop $B$-correction becomes
$B = B_0 + \Delta B_{\mathrm{KM}}(k)$.
\end{mdframed}

\paragraph{Step A4: Identify with $\mu(\Gamma_0(137))/3$.}
If $B$ evaluated at $n^* = 137$ matches $\mu(\Gamma_0(137))/3 = 46.00$,
this provides a second derivation route independent of the V$_{\mathrm{eff}}$
minimum.
Note: the match at 0.64\% is suggestive but not yet a theorem.
\begin{mdframed}[backgroundcolor=yellow!15]
\textbf{Gap A-2} \OPEN: Show that the Hecke eigenvalue structure of the
$\Theta$-partition function $Z[\Theta]$ on $\Gamma_0(137)$ determines $B$.\\
\textit{Approach}: Compute the $L$-function $L(s, f_{137})$ for a Hecke
eigenform $f_{137}$ of weight 2 on $\Gamma_0(137)$ and compare special
values $L(1, f_{137})$ to $B_{\mathrm{phenom}} \cdot \pi^{-1}$.
\end{mdframed}

\subsection{Current Status}

| Step | Status |
|------|--------|
| A1: Winding-mode action form | \PROVED\ \Lone |
| A2: Heat kernel trace | \PROVED\ \Lone |
| A3: Kac-Moody level $k$ | Gap A-1 \OPEN |
| A4: Hecke connection | Gap A-2 \OPEN |

\textbf{Blocker}: Both gaps are open. Route A is on \textbf{4-week time-box}.

%──────────────────────────────────────────────────────────────────────────────
\section{Route B: V$_{\mathrm{eff}}$ Spectral Route (PRIMARY)}
\label{sec:route_b}
%──────────────────────────────────────────────────────────────────────────────

\subsection{Central Claim}

\begin{tcolorbox}[colback=green!4!white,colframe=green!60!black,
  title=Route B Target (PRIMARY ROUTE)]
Close Gap G137-B: derive $B_{\mathrm{phenom}} \approx 46.298$ from $S[\Theta]$.
Given the proved chain (Theorems~\ref{thm:neff}--\ref{thm:b0}),
this immediately promotes $n^*(B) = 137$ from CONDITIONAL to PROVED.
\end{tcolorbox}

\subsection{Equation Chain}

\paragraph{Step B1: Full one-loop $B$ as heat-kernel sum.}
Beyond the baseline $B_0 = 8\pi$, the full one-loop coefficient is:
\[
  B \;=\; B_0 + \sum_{k=1}^{\infty} c_k \cdot a_k,
\]
where $a_k$ are Seeley-DeWitt heat-kernel coefficients for the operator
$\nabla^\dagger\nabla$ on $S^1_\psi \times M^4$, and $c_k$ are
combinatorial factors from the $N_{\mathrm{eff}} = 12$ mode count.

\paragraph{Step B2: Higher-loop correction.}
The two-loop contribution on $S^1_\psi$ is:
\[
  \Delta B^{(2)}
  \;=\;
  \frac{N_{\mathrm{eff}}}{4\pi}\cdot
  \int_0^\infty \frac{\mathrm{d}t}{t}\,e^{-m^2 t}\,K_2(t;R_\psi),
\]
where $K_2(t)$ is the two-loop heat kernel and $m$ is the winding-mode mass.
\begin{mdframed}[backgroundcolor=yellow!15]
\textbf{Gap G137-B (main blocker)} \OPEN: Evaluate $\Delta B^{(2)}$ and
show $B_0 + \Delta B^{(2)} = B_{\mathrm{phenom}} \approx 46.298$, or
equivalently $\Delta B^{(2)} \approx 21.16$.\\
\textit{Known}: $\Delta B^{(2)} > 0$ (positive loop correction).\\
\textit{Unknown}: Whether two-loop alone is sufficient, or whether
higher-loop / non-perturbative contributions are needed.
\end{mdframed}

\paragraph{Step B3: Modular bootstrap alternative.}
An alternative to loop expansion: identify $B$ with a modular-form special value.
The corroboration $\mu(\Gamma_0(137))/3 \approx B_{\mathrm{phenom}}$
(0.64\% match) suggests:
\[
  B \;\stackrel{?}{=}\; \frac{\mu(\Gamma_0(137))}{3}
    \;=\; \frac{138}{3} = 46.
\]
\begin{mdframed}[backgroundcolor=yellow!15]
\textbf{Gap G137-B-mod} \OPEN: Prove (or disprove) that the modular
measure $\mu(\Gamma_0(n^*))$ coincides with $3B$ evaluated at the prime
attractor $n^*$, using only $S[\Theta]$ as input.
\end{mdframed}

\subsection{Current Status}

| Step | Status |
|------|--------|
| B1: Shared foundation | \PROVED\ \Lone |
| B2: Two-loop evaluation | Gap G137-B \OPEN |
| B3: Modular bootstrap | Gap G137-B-mod \OPEN |

\textbf{Blocker}: Gap G137-B. Route B is PRIMARY; 4-week modular bootstrap attack.

%──────────────────────────────────────────────────────────────────────────────
\section{Route C: EW/GUT Bridge}
\label{sec:route_c}
%──────────────────────────────────────────────────────────────────────────────

\subsection{Central Claim}

\begin{tcolorbox}[colback=orange!5!white,colframe=orange!60!black,
  title=Route C Target]
If $\CC\otimes\HH$ embeds into a GUT group fixing $\sin^2\theta_W(\mathrm{GUT}) = 3/8$
(e.g.\ via SU(5) or SO(10) embedding), then renormalization-group running gives:
\[
  \sin^2\theta_W(M_Z) \;\approx\; 0.231
  \quad\Longrightarrow\quad
  \alpha^{-1}(M_Z) \;\approx\; 128.
\]
Running further to low energies gives $\alpha^{-1}(0) \approx 137.036$.
\end{tcolorbox}

\subsection{Equation Chain}

\paragraph{Step C1: Identify the GUT embedding.}
The biquaternion algebra $\BB \cong \mathrm{Mat}(2,\CC)$ has
$\mathrm{Aut}(\BB) \cong [\mathrm{GL}(2,\CC)\times\mathrm{GL}(2,\CC)]/\ZZ_2$.
Does this automorphism group embed into SU(5) or SO(10)?
\begin{mdframed}[backgroundcolor=yellow!15]
\textbf{Gap C-1} \OPEN: Determine whether $\mathrm{SU}(3)\times\mathrm{SU}(2)_L\times\mathrm{U}(1)_Y
\subset \mathrm{Aut}(\BB)$ admits an extension to $\mathrm{SU}(5)$ in which the
fundamental representation $\mathbf{5}$ of SU(5) restricts to the correct SM
representation content.
\end{mdframed}

\paragraph{Step C2: Fix $\sin^2\theta_W(\mathrm{GUT}) = 3/8$.}
In SU(5), the GUT prediction is:
\[
  \sin^2\theta_W(M_{\mathrm{GUT}}) \;=\; \frac{3}{8}.
\]
This follows from the normalization of the hypercharge generator
$Y = \mathrm{diag}(-2/3,-2/3,-2/3,1,1)$ in the $\mathbf{5}$ of SU(5).

\paragraph{Step C3: RG running to $M_Z$.}
With three SM generations and the SM $\beta$-functions:
\[
  \sin^2\theta_W(M_Z)
  \;=\;
  \frac{\sin^2\theta_W(M_{\mathrm{GUT}})}{1 + \delta_{\mathrm{RG}}}
  \;\approx\; 0.231,
\]
where $\delta_{\mathrm{RG}}$ is the one-loop RG correction from
$M_{\mathrm{GUT}}$ to $M_Z$.
This is a \textbf{standard SM calculation} conditioned on the GUT embedding.

\paragraph{Step C4: Extract $\alpha$.}
With $\sin^2\theta_W(M_Z) = 0.231$ and the SM relation
$e = g\sin\theta_W$:
\[
  \alpha(M_Z) \;=\; \frac{g^2\sin^2\theta_W}{4\pi}
             \;\approx\; \frac{1}{128},
\]
and running to $q^2 = 0$ via QED gives $\alpha^{-1}(0) \approx 137.036$.

\subsection{Current Status}

| Step | Status |
|------|--------|
| C1: GUT embedding of ℂ⊗ℍ | Gap C-1 \OPEN |
| C2: GUT prediction $\sin^2\theta_W = 3/8$ | \PROVED\ (standard SU(5)) |
| C3: RG running | \PROVED\ [STD] |
| C4: α extraction | \PROVED\ [STD] |

\textbf{Blocker}: Gap C-1 — GUT embedding not yet shown.
\textbf{Note}: Route C uses $g$ as an input from gauge sector; $\alpha$
derivation is complete only if $g$ is also determined by UBT.
This requires Gap G-strong (strong coupling from algebra) to be resolved first.

---

\textbf{Status}: Route C is the weakest of the three.
It is ACTIVE but CONDITIONAL on two separate gaps (C-1 and G-strong).
If neither closes in 4 weeks, Route C will be \textbf{parked}.

%──────────────────────────────────────────────────────────────────────────────
\section{Comparison Table}
\label{sec:comparison}
%──────────────────────────────────────────────────────────────────────────────

\begin{center}
\begin{tabular}{llllll}
\toprule
\textbf{Route} & \textbf{Proved steps} & \textbf{Open gaps} & \textbf{Gap difficulty} & \textbf{Independence} & \textbf{Priority} \\
\midrule
A (Modular-Hecke) & 2/4 & A-1, A-2 & Medium & Full (no α input) & 4-week box \\
B (V$_{\mathrm{eff}}$) PRIMARY & 3/3 + gap & G137-B & Medium-High & Full (no α input) & PRIMARY \\
C (EW/GUT) & 3/4 & C-1, G-strong & High & Conditional (needs g) & CONDITIONAL \\
\bottomrule
\end{tabular}
\end{center}

\textbf{Recommendation}: Routes A and B are the strongest paths.
Route B is PRIMARY.
Route A is the best attack on Gap G137-B via modular bootstrap.
Routes A and B are \textbf{complementary}: both target Gap G137-B from different angles.
Route C is retained as a contingency but has the most blockers.

\ifdefined\INCLUDEMODE
\else
  \end{document}
\fi
